\documentclass[11pt]{article}
\usepackage[margin=1in]{geometry}
\usepackage[T1]{fontenc}
\usepackage{lmodern,microtype,hyperref,enumitem,amsmath,amssymb,xcolor,booktabs}
\usepackage[numbers,sort&compress]{natbib}
\hypersetup{colorlinks=true,urlcolor=blue,citecolor=blue,linkcolor=blue}
\title{Q10: Why Do Models Solve Hard Problems Yet Make Simple Mistakes?}
\author{DLS 2026 --- module author: Kukla}
\date{October 2026}
\usepackage{url}
\begin{document}
\maketitle

\section*{Learning objectives}
\begin{itemize}[leftmargin=1.4em]
\item Explain ``jagged'' capability profiles and why they differ from human competence curves.
\item Trace specific failure modes to the training objective and tokenization.
\item Predict, roughly, where a model will be unexpectedly strong or weak.
\end{itemize}

\begin{quote}\small\textbf{Companion video (12:47):} \emph{Has Generative AI Already Peaked?} --- Computerphile.\\ \url{https://www.youtube.com/watch?v=dDUC-LqVrPU}\end{quote}

\section{The puzzle}
Student response \#18 captures a common shock: a model writes a correct merge-sort or explains general
relativity, then miscounts the letters in ``strawberry'' or botches $7 \times 8$. For a human, competence
is \emph{correlated}---someone who can do the hard thing can almost always do the easy thing. For LLMs it
is not. This is the \textbf{jagged frontier}: capability varies sharply and unpredictably across tasks
that look similarly easy to us \citep{dellacqua2023jagged}. Understanding why dissolves the paradox.

\section{Root cause 1: the training objective}
A model is trained to predict the next token on its corpus, not to be uniformly competent. So capability
tracks \emph{what the training distribution rewarded}, not human intuitions about difficulty. Explaining
relativity is \emph{easy} for the model because the internet contains vast, redundant, well-structured
prose about it; precise multi-digit arithmetic is \emph{hard} because correct long calculations are
comparatively rare and the model was never optimized to run an algorithm, only to continue text. ``Hard''
and ``easy'' are defined by data frequency and the objective, not by logical complexity.

\section{Root cause 2: tokenization}
Several classic ``simple'' failures are tokenization artifacts. The model does not see letters; it sees
sub-word tokens (module on Building a Transformer). ``Strawberry'' may be a few tokens with no explicit
letter inventory, so counting r's requires reconstructing spelling the model never represents directly.
Likewise, numbers are split into chunks that break positional arithmetic. These are not ``stupidity''---
they are predictable consequences of the input representation, and they vanish when the task is reframed
to expose the needed structure (spell it out, use digits with separators, or give a tool).

\section{Root cause 3: one forward pass, no built-in checking}
Without chain-of-thought or tools, the model commits to each token in a single pass with no scratchpad and
no verification step. Humans catch their own arithmetic slips by re-checking; a base model has no such loop
unless prompted to reason (module Q05) or given a calculator. So errors that a quick self-check would catch
slip through. This is why the \emph{same} model often fixes its ``dumb'' mistake the instant you ask it to
think step by step---the capability was latent; the harness was missing.

\section{A predictive heuristic}
You can anticipate jaggedness. Expect \textbf{strength} where the task is well-represented in text,
tolerant of small errors, and pattern-like. Expect \textbf{weakness} where it requires exact symbol
manipulation, long error-free chains, information below the token level (characters, digits), or
knowledge rare in the corpus. The companion code (\texttt{code/jagged\_demo.py}) illustrates the
tokenization angle with a character-counting example and shows how reframing removes the error.

\section{A gallery of predictable failures}
The three root causes let us explain a whole gallery of ``how is it this dumb'' moments without surprise.
\emph{Letter counting and spelling} fail because tokens hide characters. \emph{Multi-digit multiplication}
fails because digits are chunked and the model was never trained to execute an algorithm, only to continue
plausible text. \emph{Reversing a string} or counting words fails for the same sub-token reason.
\emph{Rare factual lookups} fail because the fact was infrequent in the corpus, even when adjacent, more
common facts are solid. \emph{Consistency over long outputs} fails because there is no global planner
enforcing coherence token to token. In every case the failure is not random---given the objective, the
tokenizer, and the single-pass constraint, it is \emph{predictable in advance}. That predictability is the
useful part: it turns ``the model is unreliable'' into ``the model is unreliable \emph{here, for this
reason}, and reliable \emph{there}.''

\section{Contrast with human error}
A human who can integrate by parts will essentially never miscount the letters in a short word; their
competence is anchored to understanding. The model's competence is anchored to \emph{statistics of text},
so the two error surfaces simply do not line up. This is why our intuitions mislead us so badly when
judging models: we unconsciously use ``can do the hard thing'' as a proxy for ``can do the easy thing,''
a proxy that is valid for people and invalid for LLMs. Recognizing that the two systems have different
competence geometries is the single most useful mental adjustment a new practitioner can make.

\section{Why this matters}
Jaggedness means you cannot infer a model's reliability on task B from its brilliance on task A. For
deployment this is central: pair models with verifiers and tools for the brittle regimes (arithmetic,
counting, exact retrieval), and do not trust ``it's smart, so it must get this right.'' The evaluation
track develops how to map a model's frontier rather than assuming a single competence level.

\section{Summary}
Models make simple mistakes because ``simple'' is a human notion; the model's competence tracks training
frequency and the next-token objective, is distorted by tokenization, and lacks a built-in checking loop.
The capability profile is jagged and must be \emph{measured}, not assumed---and brittle regimes are best
handled with reasoning prompts, tools, and verifiers.

\bibliographystyle{plainnat}
\bibliography{bibliography}
\end{document}
